\section{Transformations de Lorentz}
\subsection{Termes bilinéaires}
Dressons une liste exhaustive des grandeurs bilinéaires $\bar{\psi}f(\g^{\mu})\psi$
afin de construire des objets invariants de Lorentz – utiles pour écrire des lagrangiens
réalistes pour les particules de spin 1/2.

À cette fin, nous introduisons la matrice $\g_5$ définie comme suit :
\[
\g^5=i\g^0\g^1\g^2\g^3\ \mathds{1}
\]
Exercice : Montrez que $\g^{5\dagger}=\g^5$, $(\g^{5})^2=I$ et
$\g^5\gm+\gm\g^5=0$.\\
Dans la représentation de Dirac-Pauli,
\[
\g^5=\begin{pmatrix}0&\Id\\\Id&0\end{pmatrix},\ avec\ \g^k=\beta\alpha^k
=\begin{pmatrix} \Id & 0 \\ 0 & -\Id \end{pmatrix}
\begin{pmatrix} 0 & \s^k \\ \s^k & 0 \end{pmatrix}=
\begin{pmatrix} 0 & \s^k \\ -\s^k & 0 \end{pmatrix}
\]
puisque,
\[
\g^0\g^1 = \begin{pmatrix} 1 & 0 \\ 0 & -1 \end{pmatrix}
\begin{pmatrix} 0 & \s^k \\ -\s^k & 0 \end{pmatrix}=
\begin{pmatrix} 0 & \s^1 \\ \s^1 & 0 \end{pmatrix}
\]
\[
\g^0\g^1\g^2 = 
\begin{pmatrix} -\s^1\s^2 & 0 \\ 0 & \s^1\s^2 \end{pmatrix}
\]
\[
\g^0\g^1\g^2\g^3 = 
\begin{pmatrix} -\s^1\s^2 & 0 \\ 0 & \s^1\s^2 \end{pmatrix}
\begin{pmatrix} 0 & \s^3 \\ -\s^3 & 0 \end{pmatrix} = 
\begin{pmatrix}0&-\s^1\s^2\s^3\\-\s^1\s^2\s^3&0\end{pmatrix}
\]
\[
\ent\g^3 =
\begin{pmatrix}0&-i\s^1\s^2\s^3\\-i\s^1\s^2\s^3&0\end{pmatrix}
= \begin{pmatrix} 0_{2\times 2} & 
\begin{matrix} 1 & 0 \\ 0 & 1 \end{matrix}
 \\ 
\begin{matrix} 1 & 0 \\ 0 & 1 \end{matrix}
 & 0_{2\times 2} \end{pmatrix}
\]
comme,
\[
\s^1\s^2=\begin{pmatrix} 0 & 1 \\ 1 & 0 \end{pmatrix}
\begin{pmatrix}0&-i\\i&0\end{pmatrix}=\begin{pmatrix}i&0\\0&-i\end{pmatrix}
\]
\[
\s^1\s^2\s^3=\begin{pmatrix} i & 0 \\ 0 & -i \end{pmatrix}
\begin{pmatrix}1&0\\0&-1\end{pmatrix}=\begin{pmatrix}i&0\\0&i\end{pmatrix}
\]
Liste des grandeurs bilinéaires :

\begin{tabular}{lll}
Nom & Expression & Nombre de composantes\\
scalaire & $\overline{\psi}\psi$ & 1\\
vecteur & $\overline{\psi}\g^{\mu}\psi$ & 4\\
tenseur & $\overline{\psi}\s^{\mu\nu}\psi$ & 6\\
vecteur axial & $\overline{\psi}\g^5\g^{\mu}\psi$ & 4\\
pseudo-scalaire & $\overline{\psi}\g^5\psi$ & 1\\
 &  & \\
\end{tabular}

$\bullet$ Ordre croissant des $\g$ : $\g^5\gm=i\g^0\g^1\g^2\g^3
\gm\Leftrightarrow$ produit de 3 matrices $\g$ car $(\gm)^2=\pm\gm$


$\bullet$ ex : $\g^5\g^2\g^3=i\g^0\g^1\g^2\g^3\g^2\g^3$ est un produit de $2\g\ent$
c'est une liste exhaustive des possibilités.

\subsection{Invariants de Lorentz}

On s'attend à ce que, sous l'effet d'un boost de Lorentz,

\[
i\gm\frac{\partial}{\partial x^\mu}\psi(x)-m\psi(x)=0\quad\mapsto
\quad i\gm\frac{\partial}{\partial x'^\mu}\psi'(x')-m\psi'(x')=0
\]

\[
\text{avec la transformation de Lorentz :}\qquad x^\mu = \Lambda^{\mu.}_{.\nu} x^\nu{'}
\]
\[
\text{et la transformation quantique :}\qquad \wt{X}'^\mu = T\;\wt{X}^\mu\;T^{-1}
\]
\[
\text{de}\qquad \bra{\psi'}X'\ket{\psi'}=\bra{\psi}X\ket{\psi}
\]
\[
\bra{\psi}T^\dagger X'T\ket{\psi}= \qquad \text{« invariance de mesure »}
\]
ainsi que la transformation du champ spinoriel :
\[
\psi(x^\nu)\quad\mapsto\quad \psi'(x^\nu{'})=T\;\psi(x^\nu)
\]
Dans le référentiel F,
\[
i\g^\r\ub{\partial_\r}{\Lambda^{\mu.}_{.\r}\partial_\mu'}\psi(x)-m\;\psi(x)=0
\qquad \partial_{x^{{'}\mu}}=\Lambda^{.\nu}_{\mu.}\partial_{x^\nu}
\]
\[
\rightarrow \Lambda^{\mu.}_{.\r}\partial_\mu'=\ub{\delta^\nu_\r\partial_\nu}{\partial_\r}
\]
tandis que dans le référentiel F',
\[
iT^{-1}\g^\mu\partial_{x^{{'}\mu}}\ub{T\psi(x)}{\psi'(x')}-m\ub{T^{-1}\psi'(x')}{\psi(x)}=0
\]
Par identification,
\[
\boxed{\Lambda^{\mu.}_{.\r}\g^\r=T^{-1}\g^\mu T}\quad\ssi
\ub{\Lambda^{\mu.}_{.\r}(-\b)\g^\r}{\text{transfert classique avec }-\b}=
\ub{T\g^\mu T^{-1}(\b)}{\text{transfert quantique par définition }(\g{'}^{\mu})}
\]
Exercice :
Pour la transformation de Lorentz infinitésimale $\Lambda_{\cdot\mu}^{\nu} = 
\delta_{\mu}^{\nu} + \epsilon_{\cdot\mu}^{\nu}$.
vérifiez que la condition $T(\d^\mu_\r+\e^{\mu.}_\r)\g^\r=\g^\mu T$ est bien satisfaite par
$\boxed{T=\Id-\frac{i}{4}\;\s_{\a\b}\;\e^{\a\b}}$
\[
\tag{}  \qquad %\{}{}{}{}{} \ \ \[\] \qquad %\{} \[\]\quad\quad
\]

\[
\begin{array}{ r c l }
(\cancel{\d^\mu_\r}+\e^{\mu.}_{.\r})\g^\r-\frac{i}{4}\s_{\a\b}\,\e^{\a\b}(\d^\mu_\r+\e^\mu_\r)\g^\r
 & = & \cancel{\g^\mu}-\frac{i}{4}\g^\mu\e_{\a\b}\,\e^{\a\b} \\
\e^\mu_\r\g^\r+\frac{i}{8}[-2\g_\b\g_\a]\g_\nu\e^{\a\b}\g^{\mu\nu}
 & = & \frac{i}{8}\g^\mu[\g_\a\g_\b-\g_\b\g_\a]\e^{\a\b} \\
  & = & -\frac{i}{4}\,\g_\nu\,\g_\b\,\g_\a\,\e^{\a\b}\,g^{\mu\nu} \\
\end{array}
\]

\[
\begin{array}{ r c l }
\e^\mu_\r\g_\r-\frac{i}{4}(-)2\,\ub{g_{\b\nu}\,g^{\ob{\nu\mu}{sym.}}}
 {\d^\mu_\b}\,\g_\a\,\e^{\a\b}-\frac{i}{4}\,\g_\r\,2\,\ub{g_{\a\nu}\,g^{\nu\mu}}
 {\d^\mu_\b}\,\e^{\a\b} & = & 0 \\
\e^{\mu\r}\g_\r+\ub{\frac{1}{2}\g_\a\e^{\a\mu}}{-\frac{1}{2}\g_\a\e^{\mu\a}}
 -\frac{1}{2}\g_\b\e^{\mu\b} & = & 0 \\
\e^{\mu\r}\g_\r-\g_\a\e^{\mu\a} & = & 0 \\
\end{array}
\]
Exercice :Montrer que : $\boxed{T^{\dagger}\g^{0}=\g^{0}T^{-1}}$ d’ordre $\e^{\alpha
\beta}$.Remarque : il convient de traiter avec soin la relation
$T^{\dagger}=T^{-1}$ à $\mathcal{O}(\epsilon)$.
Cela revient à montrer que,
\[
\begin{array}{ r c l }
 T^\dagger\g^0T\g^0 & = & \Id \\ \ssi\Id & = & \big(1+\frac{1}{8}[
 \ub{\gai\gbi-\gbi\gai}{\ub{\gbid\gaid-\gaid\gbid}{-(-\gbid\gaid+\gaid\gbid)}}
]\e^{\a\b}\big)
\ub{\g_0\big(1+\frac{1}{8}[\gai\gbi-\gbi\gai]\e^{\a\b}\big)\g_0}
 {1+\frac{1}{8}[\gaid\gbid-\gbid\gaid]\e^{\a\b}} \\
\end{array}
\]
Invariants de Lorentz (infinitésimaux) :

Un premier invariant est :

\[
\begin{array}{ r c l }
 \overline{\psi'}\psi' & = & (\psi')^\dagger\g^0T\psi \\
 & = & (T\psi)^\dagger\g^0T\psi \\
 & = & \psi^\dagger T^\dagger\g^0T\psi \\
 & = & \overline{\psi}\psi \\
\end{array}
\]

Par ailleurs, on peut désormais écrire :
\[
\overline{\psi'}\gm\psi' = \overline{\psi}\ub{T^{-1}\gm T}{\Lambda^{\mu.}_{.\r}\g^\r}\psi
\]
ainsi que,
\[
\d'_\mu=\Lambda^{.\s}_{\mu.}\ds=\Lambda^{-1\s.}_{.\mu}\ds
\]
de sorte qu’un autre invariant peut être construit,
\[
\overline{\psi'}\dm'\gm\psi' = \overline{\psi}\ub{\Lambda^{-1\s.}_{.\mu}\Lambda^{\mu.}_{.\r}}{matrice\ \delta^\s_\r}\gr\ds\psi = 
\overline{\psi}\gr\dr\psi
\]
